\section{Game domain and rulesets}
\label{sec:game}

\subsection{CellGameModule}
Primary gameplay module. Owns a \symbolref{CellContext} and an enum
\symbolref{CellState}:
\begin{center}
\texttt{NORMAL} $\mid$ \texttt{EDIT} $\mid$ \texttt{EXIT}
\end{center}

\begin{statusbox}
Mode handling is an \textbf{enum + switch} inside the module, not the older
class-per-state hierarchy (archived under \pathref{archive/old-states/}).
\texttt{E} toggles EDIT/NORMAL and wakes \symbolref{SplashText}.
Save/load, pause/run/step, canvas, ruleset, and camera behavior is exposed as
registered console commands rather than artificial frame states.
Simulation rate: env \texttt{tps} $\times$ \texttt{speedFactor}, with
accumulator + step cap. Painting and simulation operate on signed world
coordinates; there is no finite or toroidal production canvas.
\end{statusbox}

\subsection{Domain console commands and persistence}
\symbolref{CellGameModule} registers domain commands through
\symbolref{CommandRegistry}; Services remains free of Game types. Save always
writes version 2 sparse files with magic/version, ruleset, camera metadata, and
sorted 16$\times$16 chunk records. Load validates temporary state before
mutation, accepts both version 2 and the prior dense format, imports legacy
cells around the origin, restores ruleset/camera metadata, and snaps the
bounded view. Native picker variants route through the platform
\symbolref{SaveLoad} API. See D-CLI1 and
\cref{sec:session-2026-08-04}.

\subsection{CellContext}
Constructs \symbolref{SparseCellGrid} and bounded \symbolref{CanvasView} from
env vars (\texttt{CanvasX}, \texttt{CanvasY}) and selects a \symbolref{RuleSet}
from a mode string
(\texttt{GAME\_OF\_LIFE}, \texttt{BRIANS\_BRAIN}, \texttt{SEEDS}, \ldots).
Mode changes from console/env call \texttt{setRuleSet} and
\texttt{CanvasView::rebuildPalette}.

\subsection{SparseCellGrid (domain) and CanvasView (view --- D-C3)}
\begin{itemize}
  \item Domain: signed 64-bit cell addresses map through floor division into
        non-background 16$\times$16 byte chunks. Negative coordinates and all
        multi-state values are preserved; no fixed chunk pool cap exists.
  \item Simulation: active chunks plus their neighbors are read through a
        temporary 18$\times$18 halo and stepped serially with stateless
        \texttt{RuleSet::nextState}; empty space is implicit and non-toroidal.
  \item View: \symbolref{CanvasView} samples the camera-visible region into a
        bounded RGB staging buffer, fades palette colors, and snaps newly
        revealed cells.
  \item GPU: one reusable RGB display texture and one screen-space
        \symbolref{GameVisual} quad; at most one dirty texture update per view
        submission.
  \item Dense \symbolref{CellGrid}/\symbolref{Canvas} remains only for
        compatibility tests and legacy import, not as a second runtime path.
\end{itemize}

\subsection{Scale walls (simulation / memory)}
\begin{statusbox}
Simulation cost follows allocated chunks and their eight-neighbor halo set,
not a fixed world rectangle. Presentation cost follows the bounded
\texttt{CanvasX}$\times$\texttt{CanvasY} view. GPU resources do not scale with
chunk count.
\end{statusbox}

\subsection{Rulesets}
Active rules (in \pathref{AllSets.h} / factory): Game of Life, Seeds, Brian's
Brain, Highlife, Day \& Night, Life Without Death, Wireworld. Stub headers
remain for Rule 90 / 184 (identity \texttt{nextState}).

\begin{lstlisting}
class RuleSet {
  // Pure transition: cell + Moore count of alive (value==0) neighbors.
  virtual unsigned char nextState(unsigned char cell,
                                  unsigned char aliveNeighbors) const;
  virtual void evalCell(const unsigned char& target,
                        unsigned char dest[3]) const; // palette colors
};
\end{lstlisting}

\textbf{Generation:}
\begin{enumerate}
  \item Collect every allocated chunk and its eight neighboring chunk
        addresses.
  \item Fill a temporary 18$\times$18 read-only halo for each target chunk.
  \item Evaluate the 16$\times$16 core and install only non-background chunks.
\end{enumerate}

\begin{decisionbox}[title={Keep rules pure}]
Rulesets stay free of rendering and input. They read/write cell buffers only
and expose \texttt{evalCell} for palette colors. That keeps them testable
(headless \texttt{TestRuleSets}).
\end{decisionbox}

\subsection{Cell encoding}
\begin{center}
\begin{tabular}{@{}ll@{}}
\toprule
\textbf{Value} & \textbf{Meaning} \\
\midrule
$0$ & Alive / Wireworld head (neighbor count treats as ``live'') \\
$1$ & Dead / Wireworld empty \\
$2$ & Multi-state (Brian's Brain dying; Wireworld tail) \\
$3$ & Wireworld conductor \\
\bottomrule
\end{tabular}
\end{center}
Historical neighbor tests used \texttt{!getCanvasPixel} (alive when zero).
Wireworld reuses head $=0$ so head-neighbor counting stays simple.
